% Related literature. States no result of this study. Every published number is quoted as
% sources/notes.md records it (sections 1-3, 9, 10) and as the citation claims in claims.jsonl
% state it. Units differ: Clemens et al. report index points (percent of the September 2023
% level), Dube log points times 100.
\section{Related literature}
\label{sec:literature}

\subsection{Estimates of AB 1228's employment effect}
\label{sec:lit-estimates}

Five studies estimate the employment effect of AB 1228 with administrative data, and they
disagree. Table~\ref{tab:published} lists the estimates closest to ours.

\citet{clemens2025did} compare California's monthly QCEW fast-food employment, indexed to
September 2023, with the rest of the United States. Between September 2023 and September 2024
it fell 2.7 percent relative to the comparison, 3.2 percent after removing pre-enactment trends,
and their eight difference-in-differences (DD) and triple-difference estimates range from $-2.28$
to $-3.93$ percent (their Table~1). They report placebo rank $p$-values, not standard errors. The
revised and renamed version \citep{clemens2026employment} extends the QCEW window to September
2025 and adds the Quarterly Workforce Indicators (QWI): its eight two-year QCEW estimates range
from $-2.02$ to $-5.03$ percent with a median of $-3.45$ percent (Table~3), and the two-year QWI
median is $-3.52$ percent (Table~7). They put the fast-food wage effect at just over 6 percent by
2025Q3 and the two-year own-wage elasticity at just over $-0.5$. We refer to these authors as
CEMN.

\citet{dube2026labor} uses quarterly QCEW data through 2025Q3, the log of employment per capita,
and as controls the 20 states whose binding fast-food minimum is the federal \$7.25. Measured as
the 2024Q2--2025Q3 average relative to 2023Q4, with 2024Q1 dropped as a transition quarter, his
event-study DD gives an employment effect of $-1.39$ percent (Ferman--Pinto standard error 0.44)
and his preferred SDID estimate is $-0.26$ percent (0.54), with wage effects of 7.50 and 6.99
percent (Table~1, Panel~A). With the 2022Q1--2023Q3 average as base (Panel~B) the two employment
estimates are $-1.28$ and $-0.02$ percent.

\citet{sosinskiy2026effects} use county QCEW data for 2022--2024, 2024Q1 as reference period and a
triple difference against full-service restaurants, and estimate an employment-rate effect of
$-0.9$ percent in 2024Q4, with a Rambachan--Roth 95 percent interval of $-2.9$ to $0.9$ percent
(Table~1). Across 1,440 synthetic-control and SDID specifications on QCEW and QWI data through
2025Q4, \citet{clemens2026synthetic} find that over 96 percent of the estimated employment effects
are negative, that the median QCEW SDID estimate is $-3.58$ percent, and that Dube's SDID
estimate lies near the 97th percentile; detrending is the single choice that moves the estimates
most, and they report that Dube's SDID omits the intercept and uses one penalty for unit and time
weights. With establishment-level YTS data through 2025, \citet{melo2026covered} find a small,
insignificant state-level effect (an inverse-hyperbolic-sine coefficient of $-0.0137$, placebo
standard error 0.0240) but significant declines among smaller chains, and no significant effect
for the 20 largest chains.

Other evidence uses other data. From Equifax payroll records of large fast-food employers,
\citet{hamdi2025wage} find that hourly wages rose 12.1 percent and establishment employment
changed by $+2.3$ percent (95 percent interval $-1.0$ to $5.5$). The first version of
Sosinskiy and Reich's study \citep{reich2024sectoral} found from job postings that pay rose 18
percent without a fall in employment, and a worker survey found large wage increases with no
effect on hours or scheduling \citep{schneider2024early}.

No causal or pre-specified estimate we found uses QCEW data for any quarter of 2026: the latest
use QCEW data through 2025Q4 \citep{clemens2026synthetic} or 2025Q3
\citep{dube2026labor,clemens2026employment}. The BLS released preliminary 2026Q1 data on August
28, 2026 \citep{bls2026cewq1}. A descriptive comparison is public: \citet{steier2026recordinglaw}
indexes California and rest-of-US 722513 employment to September 2023 through March 2026, putting
California at 97.06 and the rest of the United States at 98.99 in March 2026, and indexes average
weekly pay for 2026Q1 to 2023Q3 (114.8 against 104.3); the page makes no causal estimate. We
disclose in Section~\ref{sec:prespec} that these values were read before the design was fixed.

\begin{table}[tbp]
\centering
\caption{Published QCEW-based estimates of AB 1228's employment effect closest to this study}
\label{tab:published}
\small
\begin{tabularx}{\linewidth}{@{}>{\raggedright\arraybackslash}p{2.6cm}>{\raggedright\arraybackslash}p{2.0cm}>{\raggedright\arraybackslash}X>{\raggedright\arraybackslash}p{2.3cm}>{\raggedright\arraybackslash}p{3.0cm}@{}}
\toprule
Study & Last QCEW quarter used & Estimator and comparison & Base & Employment effect (percent) \\
\midrule
CEMN (2026), Table~3, Panel~B & 2025Q3 & Two-year fixed-horizon DD of indexed levels, rest of US; detrended & September 2023 & $-2.74$; detrended $-3.72$; median of eight $-3.45$ \\
\addlinespace
Dube (2026), Table~1, Panel~A & 2025Q3 & 2024Q2--2025Q3 average; DD and SDID on log employment per capita; 20 federal-minimum states & 2023Q4, 2024Q1 dropped & DD $-1.39$ (0.44); SDID $-0.26$ (0.54) \\
\addlinespace
Dube (2026), Table~1, Panel~B & 2025Q3 & as Panel~A & 2022Q1--2023Q3 mean & DD $-1.28$; SDID $-0.02$ \\
\addlinespace
Sosinskiy and Reich (2026), Table~1 & 2024Q4 & County triple difference against 722511, employment rate & 2024Q1 & $-0.9$ [$-2.9$, $0.9$] in 2024Q4 \\
\addlinespace
Clemens et al.\ (2026, w35931) & 2025Q4 & 1,440 synthetic-control and SDID specifications & 2023Q4 weights; 2023Q3 for fixed horizons & median QCEW SDID $-3.58$ \\
\bottomrule
\end{tabularx}

\medskip
\begin{minipage}{\linewidth}\footnotesize
\emph{Notes:} As recorded from each source (sources/notes.md). Parentheses: Ferman--Pinto
placebo standard errors; brackets: Rambachan--Roth 95 percent interval. CEMN report index points
(percent of the September 2023 level); Dube reports log points times 100. The estimands differ
from ours and from each other in horizon, base and outcome normalisation.
\end{minipage}
\end{table}

\subsection{Why the estimates differ, and what is already mapped}
\label{sec:lit-map}

The disagreement has been decomposed. Starting from \citet{clemens2025did}, \citet[Table~D1]{dube2026labor}
changes one specification choice at a time and attributes most of the gap to detrending,
normalisation by population, multi-quarter pre- and post-period averaging windows (one step that
replaces the September 2023 base and single post month with the 2022Q1--2023Q3 and 2024Q2--2025Q3
averages) and synthetic-control weighting, and finds seasonal adjustment and the choice of
control states nearly irrelevant; on QWI data the switch of data source alone accounts for
roughly half of the gap. \citet[Table~A1]{clemens2026synthetic} report the median estimate by
each choice across their 1,440 specifications, and \citet[Appendix~A]{sosinskiy2026effects} list
nine differences between the studies and test four. We cite these maps and do not repeat them as
findings. We use them instead to fix the set of candidate specifications from which a
pre-specified rule picks our primary estimator (Section~\ref{sec:select}).

The base period and the anticipation window are also already varied. Clemens and coauthors
measure effects from the September 2023 enactment and argue that firms responded in
anticipation; Sosinskiy and Reich use the quarter before implementation as reference and find
no anticipation effect; Dube uses 2023Q4 as base and drops 2024Q1
\citep{clemens2026employment,sosinskiy2026effects,dube2026labor}. Re-estimating with a September
2023 reference, \citet[Table~A3]{sosinskiy2026effects} estimate employment-rate effects of $-0.2$
percent in 2023Q4 (95 percent interval $-1.9$ to $1.7$) and $-0.6$ percent in 2024Q1 ($-4.5$ to
$2.7$) in their preferred specification, and $-0.5$ and $-1.4$ percent in the specification closest
to Clemens and coauthors. \citet[footnote~5]{clemens2026synthetic} fit their synthetic controls with
Dube's base and transition quarter and report that the distribution of estimates is slightly
more negative, but not meaningfully changed, with a 2023Q3 baseline and no transition quarter.
Re-estimated on Dube's 2023Q4 base, the one-year QWI estimates of \citet[Table~A3]{clemens2026employment}
have a median of $-1.43$ percent, against $-1.38$ percent on their 2023Q3 base. Prices began to
rise between enactment and implementation \citep{clemens2026prices}, which is one reason we hold
the two quarters between signing and the effective date out of both the base and the post window
(Section~\ref{sec:timing}).

\subsection{Minimum wages and restaurant employment}

Fast-food restaurants have been the reference case for one-state minimum-wage comparisons since
\citet{card1994minimum}, who found no indication that New Jersey's 1992 increase reduced
employment in surveyed restaurants relative to eastern Pennsylvania. Across 138 state increases,
\citet{cengiz2019effect} find the number of low-wage jobs essentially unchanged five years after an
increase. For restaurants, \citet{dube2024restaurant} argue that earlier disemployment estimates
were confounded by pre-trends amplified by two-way fixed effects, while
\citet{neumark2026modern} re-estimate modern event-study designs on the same data, find negative
and significant effects, more strongly in restaurants, and argue that null results depend on
choices of variables, events, samples and weighting. The last point is one reason our weighting,
pools and outcome normalisation are fixed by a rule rather than chosen. Closest in data to ours,
\citet{wiltshire2026minimum} use county QCEW 722513 data and stacked synthetic controls for
California and New York counties and find pay growth and no disemployment during the run-up to
\$15. \citet{thornberg2025jumping} documents a large benchmark revision of California's CES
payroll estimates for limited-service restaurants; we use the QCEW, the benchmark itself, and not
the CES.

\subsection{Inference and estimator choice with one treated state}
\label{sec:lit-methods}

With one treated state, cluster-robust standard errors are not valid, as the QCEW studies above
note; they use placebo rank $p$-values \citep{clemens2025did}, Ferman--Pinto standard errors
\citep{dube2026labor} and Rambachan--Roth bounds \citep{sosinskiy2026effects} instead. \citet{conley2011inference}
estimate the distribution of the DD estimator from control-group residuals, assuming errors
independent and identically distributed across groups. \citet{ferman2019inference} show that
variation in group size makes the aggregate DD model heteroskedastic, so that such tests misstate
their size, and propose a correction valid with one treated group as the number of controls grows,
under the assumption that errors differ across groups only by a scale that depends on observables
such as size; in their simulations it performs reasonably with 25 groups. Dube's implementation
scales the placebo estimates by population. The placebo variance estimator of SDID reassigns
treatment to control units and relies on homoskedasticity across units, which
\citet{arkhangelsky2021synthetic} describe as effectively necessary for inference with one treated
unit. \citet{abadie2010synthetic} introduce placebo inference for synthetic control and suggest
in-time placebos when comparison units are few. \citet{chernozhukov2021exact} permute residuals over
time rather than units, which requires many pre-treatment periods. Our inference
(Section~\ref{sec:inference}) builds on the placebo approach and replaces its homoskedasticity
assumption with a standard error that is conservative for the largest unit.

On the choice of estimator, \citet{abadie2021using} recommends computing, and preregistering,
synthetic-control weights before post-treatment outcomes are seen, as a safeguard against
specification searches, and using backdated interventions as in-time placebo tests.
\citet{ferman2020cherry} show that with a finite pre-period there is substantial room for
specification searching across synthetic-control specifications, even among those with good
pre-treatment fit, and recommend reporting several. Demeaning outcomes by their pre-period means
relaxes the no-intercept restriction of synthetic control \citep{doudchenko2016balancing}, and
\citet{ferman2021synthetic} show that the demeaned estimator is asymptotically unbiased under
conditions under which DD is, while the original can be biased with imperfect pre-treatment fit.
Our equivalence test is the two one-sided tests procedure, which declares equivalence within
pre-specified bounds when the 90 percent confidence interval lies inside them
\citep{schuirmann1987comparison,lakens2017equivalence}.

\subsection{What this study adds}

This is a pre-specified, not blind, update of a contested question, and its novelty is modest. The
estimand, the base period, a rule that picks the primary estimator from data ending in 2019, an
inference procedure built for one treated state, and the decision rules were fixed and reviewed
before any outcome after 2023Q3 was computed. The employment hypothesis adds an equivalence test
against a bound of 2 percent that separates the two published readings; the year-two comparison is
read next to the yearly wage effects, because the fast-food premium over the general minimum fell;
and the data add one preliminary quarter, 2026Q1, to the latest published window. The
decompositions of the disagreement and the published anticipation and base-period checks are
cited, not repeated.
